% !TEX root = ../main.tex
%%%%%%%%%%%%%%%%%%%%%%%%%%%%%%%%%%%%%%%%%%%%%%%%%%%%%%%%%%%%%%%%%%%%%%%
% Conclusion
%%%%%%%%%%%%%%%%%%%%%%%%%%%%%%%%%%%%%%%%%%%%%%%%%%%%%%%%%%%%%%%%%%%%%%%

\section{Conclusion}
\label{sec:conclusion}

We have built a stock-flow-consistent agent-based economy in which investment
is endogenous, capital accumulates, labour is hired endogenously, and the wage
responds to unemployment --- and in which, consequently, the Kaleckian question
of whether retention expands or depresses output is an outcome rather than an
assumption.

The first result is that the outcome has a shape rather than a label. The
response $Y(\rho)$ is U-shaped, and the statement that survives every choice of
estimator is about the turn rather than the sign: \emph{retention depresses
output while $\rho$ lies below a $\sigma$-dependent threshold $\rhostar$ and
expands it beyond, $\rhostar$ rises with $\sigma$, and the empirically anchored
retention rate lies below $\rhostar$ for every $\sigma \geq 0.5$}. Summarised
through one slope that structure appears as a sign frontier at $\sigstar
\approx 0.65$, which rises when wages are made flexible --- the paradox of costs
dominating substitution --- and is invariant to expectation dynamics. We use
``wage-led'' and ``profit-led'' in that reduced sense throughout, as local
descriptions of one side of a turn, and not as a verdict on the economy.

The second is that the region in which retention depresses output is not a
property of the technology alone. A balanced-budget unemployment benefit at a
replacement rate inside the OECD band \emph{eliminates} it, making retention
expansionary at every elasticity tested, by closing the leak through which the
unemployed drain the circular flow. A sign reported without its institutional
setting is incomplete, and the same instrument that closes the leak is the one
that, applied to the collapse, makes it worse.

The third is the finding we would carry forward, and it was not the one we set
out to measure. It is the
wage$\,\to\,$unemployment$\,\to\,$capital-erosion channel, which emerged
unlooked-for from the wage curve, killed three consecutive stabilisation
hypotheses, and proved immune to every demand-side instrument applied to it.
Demand policy works where demand binds. Where the wage bill drives an
oscillation that erodes the capital stock, more demand is not a stabiliser ---
and in the sharpest case measured here, a saturated fiscal instrument financed
by taxing workers to pay workers simply amplifies the cycle that kills the
economy.

Two boundaries are placed on all three. Marginalised over the full defensible
parameter space the conditional structure largely dissolves: $\sigma$ explains
about 2\% of the variance and survival is governed almost entirely by the
depreciation rate, at whose empirically implied value the model does not exist
at all --- which, read with its inability to match $I/Y$ and $K/Y$ jointly,
identifies one structural culprit in a capital block without trend growth. And
the sign itself proved partly an artefact of how we measured it: a two-point
chord across a U-shaped response is not the whole-support slope that defined
the frontier, and replacing it moves the headline probability from $0.095$ to
$0.026$ while leaving the variance decomposition intact. Neither boundary was
anticipated when the quantities were chosen, which is the argument for
choosing them after the shape of the response is known rather than before.
